% 2010 年天津卷理科第 4 题：程序框图，输出 s = −7 时判断框内的条件。
% 来源 PDF：普通高考/2010/2010天津理.pdf，SHA-256 340aba76699a12f330637f9c682abbc893e03e0e91788287934c0a7e0a3cace9。
% 原页：第 1 页（A3 六栏）第 1 栏中部，书页码 635。
% 重绘依据：开始 → i=1 → s=2 → 判断；是（左）→ s=s−i → i=i+2 → 回到判断前的竖线；
%   否（下）→ 输出 s → 结束。原页判断框内为空。
\begin{tikzpicture}[semithick, >=Stealth,
  every node/.style={inner sep=2pt, font=\small}, node distance=5mm]
  \node (start) [draw, rounded corners=4pt, minimum width=2cm, minimum height=0.55cm] {开始};
  \node (i1) [draw, minimum width=1.6cm, minimum height=0.55cm, below=of start] {$i=1$};
  \node (s2) [draw, minimum width=1.6cm, minimum height=0.55cm, below=of i1] {$s=2$};
  \node (cond) [draw, diamond, aspect=2.2, minimum width=3cm,
    minimum height=1.1cm, below=16mm of s2] {};
  \node (out) [draw, trapezium, trapezium left angle=72, trapezium right angle=108,
    minimum width=2.4cm, minimum height=0.55cm, below=9mm of cond] {输出 $s$};
  \node (end) [draw, rounded corners=4pt, minimum width=2cm, minimum height=0.55cm,
    below=of out] {结束};
  \node (ss) [draw, minimum width=2.6cm, minimum height=0.55cm,
    left=13mm of cond] {$s=s-i$};
  \node (ii) [draw, minimum width=2.6cm, minimum height=0.55cm,
    above=6mm of ss] {$i=i+2$};

  \draw[->] (start) -- (i1);
  \draw[->] (i1) -- (s2);
  \draw[->] (s2) -- (cond);
  \draw[->] (cond) -- node[right] {否} (out);
  \draw[->] (out) -- (end);
  \draw[->] (cond.west) -- node[above, pos=0.72] {是} (ss.east);
  \draw[->] (ss.north) -- (ii.south);
  \draw (ii.north) -- ++(0,0.35) -| ($(s2.south)!0.5!(cond.north)$);
\end{tikzpicture}
